% Source: https://artofproblemsolving.com/community/c3h2555183_what_is_compound_interest
% Stable identity: aops-topic-2555183-q01
\textbf{(a)} The annual rate is split into four equal parts: at the end of every $3$ months, interest of $\dfrac{5\%}{4}=1.25\%$ of the \emph{current} balance (including interest already added) is added. So each quarter the balance is multiplied by $1.0125$, and after one year it is
\[1000\cdot1.0125^4=1000\cdot1.05094533\ldots\approx\boxed{\$1050.95}.\]
(Here $1.0125^2=1.02515625$ and $1.02515625^2=1.0509453369\ldots$.)

\textbf{(b)} The balance grows by the factor $1.0509453\ldots$ in a year, which is the same as simple annual interest at about $\boxed{5.09\%}$ (more precisely $5.0945\%$). Compounding makes the effective rate a little higher than $5\%$ because interest is earned on earlier interest.
